% =====================================================================
% 2 Contexto institucional e legal
% Fontes: paper/secoes/02-background.tex; docs/legislacao/README.md.
% =====================================================================

\chapter{Contexto institucional e legal}
\label{cap:contexto}

\section{A Lei Zé Maria do Tomé}
\label{sec:lei}

A Lei Estadual nº 16.820 foi sancionada em 8 de janeiro de 2019 e entrou em
vigor em 9 de janeiro. Ela acrescenta o art. 28-B à Lei Estadual nº
12.228/1993, a norma estadual sobre agrotóxicos, com a seguinte redação:

\begin{citacao}
Art. 28-B. É vedada a pulverização aérea de agrotóxicos na agricultura no
Estado do Ceará.

§ 1º A infração ao art. 1º sujeita o infrator ao pagamento de multa de 15 mil
(quinze mil) UFIRCEs.

§ 2º Fica proibida a incorporação de mecanismos de controle vetorial por meio
de dispersão por aeronave em todo o Estado do Ceará, inclusive para os casos de
doenças causadas por vírus.
\end{citacao}

Dois traços da redação organizam o desenho empírico. A proibição é de método:
nenhum princípio ativo é nomeado. E ela não se restringe à agricultura, porque
o § 2º alcança a dispersão aérea para controle de vetores. Um município sem
agricultura pulverizada, mas com controle aéreo de endemias, também é alcançado
pela norma.

A Tabela~\ref{tab:cronologia} reúne a cronologia. Ela distingue três marcos de
antecipação: a notícia (apresentação do projeto de lei), a certeza (aprovação
unânime) e a obrigação (vigência). A distinção tem consequência empírica,
discutida na Seção~\ref{sec:dose}: a janela em que se mede a exposição anterior
à lei começa depois do primeiro marco.

\begin{table}[htb]
\centering
\caption{Cronologia legal do tratamento}
\label{tab:cronologia}
\small
\begin{tabular}{lp{11cm}}
\toprule
Data & Evento \\
\midrule
09/12/1993 & Lei Estadual nº 12.228/1993, norma estadual sobre agrotóxicos. \\
20/11/2009 & Limoeiro do Norte proíbe a pulverização aérea pela Lei Municipal
             nº 1.478/2009. \\
26/05/2010 & A Lei Municipal nº 1.511/2010, de política ambiental, revoga a
             proibição de Limoeiro do Norte, um mês após o assassinato de Zé
             Maria do Tomé. \\
24/02/2015 & Apresentação do Projeto de Lei nº 18/2015 (marco de notícia). \\
18/12/2018 & Aprovação por unanimidade na Assembleia Legislativa (marco de
             certeza). \\
09/01/2019 & Vigência da Lei nº 16.820/2019 (marco de obrigação). \\
mai/2023   & O STF julga a lei constitucional (ADI 6137). \\
19/12/2024 & A Lei Estadual nº 19.135/2024 abre exceção para aeronaves
             remotamente pilotadas. \\
2025       & Nova contestação no STF (ADI 7794), pendente. \\
\bottomrule
\end{tabular}
\legend{Fonte: elaboração própria, a partir dos textos oficiais das Leis
Estaduais nº 12.228/1993 e nº 16.820/2019 e do acervo da Câmara Municipal de
Limoeiro do Norte. A revogação de 2010 está documentada por fonte secundária
(imprensa); o texto integral da Lei nº 1.511/2010 ainda será conferido.}
\end{table}

A lei de 2024 mantém a vedação, salvo quando a aplicação é feita por aeronave
remotamente pilotada, sob condições de altura, velocidade do vento e distância
de escolas, hospitais e áreas protegidas. A nova redação não reproduz o § 2º de
2019. Do ponto de vista de instrumento, o Ceará não revogou a proibição:
substituiu a cota zero por regulação de desempenho com zona-tampão. Para o
Ensaio 1, a consequência é encerrar a janela de análise em 19 de dezembro de
2024, de modo que o tratamento avaliado seja a proibição tal como vigorou entre
2019 e 2024.

\section{A Chapada do Apodi e a química do caso}
\label{sec:chapada}

A Chapada do Apodi, no Baixo Jaguaribe, concentra a fruticultura irrigada de
exportação do Ceará, sobretudo banana e melão, em perímetros de irrigação sobre
solo cárstico. É a região a que a lei responde.

A justificativa do Projeto de Lei nº 18/2015 reproduz uma tabela de análises de
água atribuída ao Dossiê ABRASCO, referente a 23 pontos de coleta na Chapada em
2009 (Tabela~\ref{tab:abrasco}). O perfil é de fungicida e inseticida sobre
fruticultura, não de herbicida sobre grão transgênico.

\begin{table}[htb]
\centering
\caption{Princípios ativos detectados em água na Chapada do Apodi, 2009}
\label{tab:abrasco}
\small
\begin{tabular}{llc}
\toprule
Princípio ativo & Classe & Amostras positivas \\
\midrule
Procimidona  & Fungicida                    & 23 de 23 \\
Carbaril     & Inseticida (carbamato)       & 23 de 23 \\
Carbofurano  & Inseticida (carbamato)       & 18 de 23 \\
Fenitrotiona & Inseticida (organofosforado) & 16 de 23 \\
Glifosato    & Herbicida                    & 4 de 23 \\
\bottomrule
\end{tabular}
\legend{Fonte: tabela reproduzida na justificativa do Projeto de Lei nº
18/2015, atribuída ao Dossiê ABRASCO. A fonte primária do Dossiê ainda será
conferida.}
\end{table}

Disso decorrem três consequências para o projeto. Os desenhos construídos sobre
soja tolerante a herbicida transferem na forma, isto é, na ideia de usar
aptidão agroclimática como fonte de variação, e não no mecanismo. A lista de
culturas usada na literatura para construir o índice de aptidão, que é de
culturas transgênicas, precisa ser adaptada, porque banana não é transgênica. E
carbamatos e organofosforados são inibidores de colinesterase, com quadro agudo
característico, o que torna a notificação de intoxicação aguda um canal
informativo.

\section{O mecanismo: deriva e substituição}
\label{sec:deriva}

O que distingue a aplicação aérea da terrestre é a deriva. A exposição
resultante é difusa e alcança a população, não apenas o trabalhador. Dois
canais operam em paralelo: o atmosférico, por deposição direta sobre quem está
a favor do vento, e o hídrico, por escoamento e percolação até a água de
consumo.

Há um terceiro canal, de sinal oposto. Proibido o avião, a aplicação não
desaparece: migra para trator e pulverizador costal. Isso afasta o produto da
população e o aproxima do aplicador. A proibição pode reduzir a exposição
populacional e aumentar a intoxicação aguda ocupacional, e é por isso que o
projeto mede esse canal em separado (hipóteses H3 e H4).

\section{Fiscalização e registros administrativos}
\label{sec:fiscalizacao}

A fiscalização do art. 28-B cabe, pela Lei nº 12.228/1993, à Superintendência
Estadual do Meio Ambiente, à Secretaria da Agricultura e à Secretaria da Saúde.
A extensão efetiva da fiscalização não é conhecida. Se a proibição não foi
fiscalizada, a exposição anterior mede intenção de aplicação, e o parâmetro
estimado é um efeito de intenção de tratar.

Três registros administrativos podem dizer quem aplicava por aeronave antes da
lei, e nenhum deles foi obtido até aqui. O primeiro é o cadastro estadual de
empresas prestadoras de serviço (art. 8º da Lei nº 12.228/1993), a cargo da
SEMACE. O segundo é o receituário agronômico, que registra a modalidade de
aplicação (Decreto Federal nº 4.074/2002, art. 66), a cargo da ADAGRI. O
terceiro são os relatórios mensais que os operadores aeroagrícolas devem
entregar à Superintendência Federal de Agricultura do estado em que atuaram,
por força da Instrução Normativa MAPA nº 2/2008. As minutas dos três pedidos de
acesso à informação estão redigidas, e o protocolo é a primeira atividade do
plano de trabalho (Capítulo~\ref{cap:plano}).
